\documentclass[12pt]{article}
\usepackage[utf8]{inputenc}
\usepackage[T1]{fontenc}
\usepackage{hyperref}
\title{Tofanma Language Resources for K-12 Curriculum Development}
\author{Compiled by Claude Sonnet 5 (claude-sonnet-5) for Andr\'as Kornai}
\date{}
\begin{document}
\maketitle

\section{Introduction}
Tofanma (ISO 639-3 tlg) is a poorly documented Papuan language of Keerom Regency, Senggi District, Papua, Indonesia. Wikipedia gives about 250 native speakers (2005). It forms the small Namla--Tofanma family with Namla. Scholars disagree on the family's placement: Usher (2020) puts it in West Pauwasi, while Foley (2018) treats it as independent.

\section{Primary Educational Resources}
The searches found no primers, textbooks or school materials for Tofanma.

\section{Linguistic and Reference Materials}
\begin{itemize}
\item Overview: \url{https://en.wikipedia.org/wiki/Tofanma_language}
\item Family article: \url{https://en.wikipedia.org/wiki/Namla\%E2\%80\%93Tofanma_languages}
\item Glottolog entry: \url{https://glottolog.org/resource/languoid/id/tofa1246}
\item New Guinea World family page: \url{https://newguineaworld.linguistik.uzh.ch/families/pauwasi-river/west-pauwasi-river/namla-tofanma/start}
\end{itemize}
These pages are summaries. The search showed no full grammar or dictionary.

\section{Religious and Community Resources}
Joshua Project reports that the Tofanma have no gospel recordings, films or videos in their language: \url{https://joshuaproject.net/index.php/people_groups/15527}.

\section{Digital and Audio Resources}
None were found. The Endangered Languages Project has a catalogue entry: \url{https://www.endangeredlanguages.com/lang/2795}.

\section{Curriculum Development Context}
Almost nothing is documented. Any curriculum would need new fieldwork and documentation with the speech community. Namla, the neighbouring language, is close to extinction and is being replaced by Tofanma and Papuan Malay.
\end{document}